% PEIT Innovation Showcase extended abstract
\documentclass[10pt,a4paper]{article}

\usepackage[margin=1.7cm]{geometry}
\usepackage[T1]{fontenc}
\usepackage[utf8]{inputenc}
\usepackage{lmodern}
\usepackage{microtype}
\usepackage{booktabs}
\usepackage{array}
\usepackage{enumitem}
\usepackage{xcolor}
\usepackage{hyperref}

\hypersetup{colorlinks=true, urlcolor=blue, linkcolor=black}
\setlist{nosep,leftmargin=*}
\setlength{\parindent}{0pt}
\setlength{\parskip}{3pt}

% The article class keeps the 2--4 page format; \chapter is mapped to a
% compact heading so chapters do not force page breaks.
\newcommand{\chapter}[1]{\section{#1}}

\title{PEIT: Personal Economic Inflation Tracker\\
\large Decentralizing Economic Truth Through Crowdsourced Purchase Analytics}
\author{Ken Tanaka Hernández$^{1}$ \and Co-author Name$^{2}$\\
\small $^{1}$HAIKU, Guadalajara, Mexico \quad $^{2}$HAIKU, Guadalajara, Mexico\\
\small \href{mailto:kentanaka3@outlook.com}{kentanaka3@outlook.com}}
\date{}

\begin{document}
\maketitle

\begin{abstract}
Personal inflation is often discussed using population-level indices that may
not represent an individual's consumption pattern. PEIT (Personal Economic
Inflation Tracker) is a community-driven financial ecosystem designed to shift
the power of macroeconomic tracking from centralized governments directly to
individual consumers.

Traditional economic metrics, such as the official Consumer Price Index (CPI),
rely on a generalized, rigid "basket of goods and services." This macro
approach inherently smooths out volatile price surges and fails to accurately
reflect diverse lifestyles, cultural nuances, or specific dietary choices. By
treating every consumer as an independent data node, this platform captures the
real-world, localized price movements of goods and services at the point of
sale.

Through standard mobile camera snapshots, the system instantly transforms
static paper and digital receipts into dynamic, structured economic data. This
data is split to achieve two simultaneous goals: maintaining a private ledger
for personal lifestyle analysis and contributing to an open-access,
crowdsourced price map.

In its current form, PEIT is a research-oriented prototype that records monthly
expenditure lines from manually transcribed receipts, separating price change
from changing purchasing behaviour. The prototype preserves the original
receipt wording, quantity, unit, signed financial amount, price components,
event type, and source-file provenance. It validates monthly CSV inputs,
appends an auditable JSON-lines import record, and produces two distinct
indicators: an exact-normalized-item fixed-basket price index and a
changing-basket paid-expenditure index. It also reports item-match coverage and
month-over-month change, making incompleteness visible rather than silently
inventing observations. The current demonstrator is intentionally modest: its
price aggregation is an unweighted arithmetic mean of matched price relatives,
and category fallback, substitution review, uncertainty intervals,
multi-currency conversion, and forecasting remain planned extensions. This
separation creates a transparent foundation for testing alternative basket
definitions and weight systems, including personal preferred weights and
historically recognized weights. During a live demonstration, a synthetic MXN
receipt file can be validated, audited, and compared across two months in a
reproducible command-line workflow. The intended outcome is not an official
inflation statistic, but an inspectable personal measurement system that can
support budgeting, research into consumption behaviour, and future user
interaction evaluation.
\end{abstract}

\textbf{Keywords:} personal inflation; expenditure tracking; price index;
reproducible analysis; receipt data; local-first software; consumption behaviour

\chapter{Background and motivation}

Inflation is experienced through the prices and quantities that an individual
actually encounters. A population-wide consumer price index is valuable for
macroeconomic comparison, but it does not answer a personal question such as:
\emph{How has the cost of the goods and services I usually consume changed?}
Moreover, a larger monthly bill can result from higher prices, buying more,
substitution, or a one-off purchase. Treating all of these effects as one
number obscures the underlying mechanism.

PEIT addresses this need with two measurement pairs. The \emph{offer} side
describes observed price terms, while the \emph{demand} side describes
observed purchase behaviour. The first milestone targets monthly manual
receipt ingestion for one personal user, with local-only storage and explicit
human review. The project therefore prioritizes traceability and honest
uncertainty over automation.

\begin{table}[h]
\centering
\caption{The two measurements kept separate by PEIT.}
\small
\begin{tabular}{p{0.22\linewidth}p{0.32\linewidth}p{0.35\linewidth}}
\toprule
Measurement & Main question & Typical drivers of change\\
\midrule
Price / offer & What does a comparable item cost? & shelf price, tax, fee, discount\\
Expenditure / demand & What did the user spend? & price, quantity, mix, one-off events\\
\bottomrule
\end{tabular}
\end{table}

\chapter{Scientific and technological foundations}

The prototype treats a receipt line as an observation and keeps derived
quantities separate from it. Its signed ledger convention is that expenditure
is negative, income is positive, discounts are positive offsets, and fees,
taxes, markups, surcharges, and deposits are negative adjustments. The schema
also requires quantity and unit, enabling future comparison per kilogram,
litre, piece, hour, token, or kWh. MXN is the default example currency, while
the currency code is stored per row.

For a normalized item $i$, let $p_{i,t}$ be its quantity-weighted average unit
price in month $t$:
\[
 p_{i,t}=\frac{\sum_{k\in\mathrm{purchases}(i,t)}q_{i,t,k}\,p_{i,t,k}}
              {\sum_{k\in\mathrm{purchases}(i,t)}q_{i,t,k}}.
\]
The price relative between a base month $0$ and current month $t$ is
\[
 r_{i,t}=\frac{p_{i,t}}{p_{i,0}},
 \qquad
 I^{\mathrm{price}}_t=100\,\frac{1}{|M_t|}\sum_{i\in M_t}r_{i,t},
\]
where $M_t$ is the set of exact normalized-item matches. The prototype uses
shelf unit prices for this demonstrator and reports the coverage
\[
 C_t=100\,\frac{|M_t|}{|B_0|}.
\]
The changing-basket expenditure index is
\[
 I^{\mathrm{exp}}_t=100\,\frac{\sum_{j\in t}|a_j|}
 {\sum_{j\in 0}|a_j|},
\]
where $a_j$ is the signed paid total and absolute values express positive
consumption cost in the display index. The first available month is set to
100.0. This is a prototype index, not yet a named weighted Laspeyres,
Paasche, Fisher, or chained index.

\paragraph{Worked example.} If one litre of milk costs MXN~2.50 in the base
month and MXN~3.00 later, the price index is $100(3.00/2.50)=120.0$. If two
litres are purchased later for MXN~6.00, the expenditure index is
$100(6.00/2.50)=240.0$; quantity change explains the difference.

\chapter{System architecture and design methodology}

The current architecture is a small Python package with a CSV schema, tests,
and a command-line interface. The data flow is:

\begin{verbatim}
manual receipt transcription
          |
          v
    monthly CSV --> schema validation --> append-only JSONL audit record
          |                                   |
          +--> exact item matching -----------+
                        |                     |
                        v                     v
                  price index         expenditure index
                        \                     /
                        coverage and MoM report
\end{verbatim}

The design follows four principles: (1) preserve source observations;
(2) label analyst-controlled normalization; (3) keep price and expenditure
measurements separate; and (4) make missing history and incomplete matching
visible. Transfers, refunds, and income remain typed ledger rows but are
excluded from the current purchase indices.

\begin{table}[h]
\centering
\caption{Prototype components and evidence produced.}
\small
\begin{tabular}{p{0.25\linewidth}p{0.38\linewidth}p{0.25\linewidth}}
\toprule
Component & Responsibility & Evidence\\
\midrule
CSV reader & preserve and validate receipt rows & row-level issues\\
Audit append & record import date and source file & JSONL ledger\\
Index calculator & compare exact keys and paid totals & JSON or text report\\
Focused tests & check validation and index separation & assertions\\
\bottomrule
\end{tabular}
\end{table}

\chapter{Prototype functionality and demonstration scenario}

The demonstrator uses synthetic data only. A monthly CSV contains fields such
as \texttt{item\_raw}, \texttt{item\_normalized}, \texttt{category},
\texttt{quantity}, \texttt{unit}, \texttt{shelf\_unit\_price},
\texttt{discount}, \texttt{fee}, \texttt{tax}, \texttt{deposit},
\texttt{paid\_total}, \texttt{event\_type}, and \texttt{source\_file}.
For example, a one-litre milk line can be recorded with a shelf price of
MXN~2.50 and a paid total of \$-2.50. A later line at MXN~3.00 yields a
fixed-basket price index of 120.0 for that item. If the later purchase also
contains twice the quantity, the expenditure index can rise to 240.0; the two
results communicate different facts.

\paragraph{Input example.} A synthetic row may contain
\texttt{milk 1 l}, quantity~1, unit~\texttt{L}, shelf unit price~2.50,
currency~\texttt{MXN}, discount~0, and paid total~\$-2.50. The raw wording is
retained while the normalized key enables matching.

A live sequence is:

\begin{enumerate}
  \item inspect the synthetic CSV and its schema;
  \item run the CLI validator and display row-level errors, if any;
  \item append the import and source identity to an audit JSONL file;
  \item calculate the two monthly indices, coverage, and price-index MoM;
  \item change one value or item key and show how the result changes.
\end{enumerate}

Representative Unix commands are:

\begin{verbatim}
PYTHONPATH=src python -m peit.cli config/receipt_schema.csv --json
PYTHONPATH=src python -m peit.cli config/receipt_schema.csv --audit /tmp/peit-audit.jsonl
\end{verbatim}

The exhibition package should include a short screen recording, a synthetic
two-month CSV, the exact command transcript, and a README explaining that
positive display costs are derived from signed ledger values.

\begin{description}
  \item[Visitor task 1:] identify raw and normalized item fields.
  \item[Visitor task 2:] correct a malformed quantity or date.
  \item[Visitor task 3:] compare price change with expenditure change.
  \item[Visitor task 4:] explain why incomplete matching lowers coverage.
\end{description}

\chapter{Novelty, contribution, and intended applications}

PEIT's contribution is methodological as well as technical: it makes the
personal inflation question explicit as a separation between price conditions,
purchase behaviour, and total expenditure. The audit trail and status labels
also prevent a normalized label or a forecast from being mistaken for a raw
observation. Potential applications include personal budgeting, reflective
consumption analysis, educational demonstrations of index construction, and
research on individualized expenditure patterns. The local-first design is
particularly suitable for sensitive financial information.

\begin{itemize}
  \item \textbf{Personal use:} understand recurring costs and unusual bills.
  \item \textbf{Education:} demonstrate how basket definitions affect indices.
  \item \textbf{Research:} study individualized price and consumption patterns.
  \item \textbf{Technology transfer:} provide a privacy-preserving foundation.
\end{itemize}

\chapter{Current development status}

The project is a working command-line proof of concept. CSV validation,
signed receipt fields, JSON-lines provenance, exact normalized-item matching,
fixed-basket price calculation, changing-basket expenditure calculation,
coverage, and tests are implemented. The web interface is intentionally
deferred. The following capabilities are specified but not yet implemented:
category-level fallback, human-reviewed equivalence, unit-conversion tables,
category contributions, uncertainty quantification, time-based currency
conversion, and multi-horizon forecasts. Results from the current prototype
must therefore be presented as derived personal indicators with limited
coverage, not as official inflation or validated forecasts.

\begin{table}[h]
\centering
\caption{Implemented and planned capabilities at the exhibition stage.}
\small
\begin{tabular}{p{0.38\linewidth}p{0.18\linewidth}p{0.32\linewidth}}
\toprule
Capability & Status & Presentation wording\\
\midrule
CSV validation and provenance & Implemented & working prototype\\
Exact-key price and expenditure indices & Implemented & demonstrable\\
Category fallback and substitutions & Planned & future work\\
Weights and uncertainty intervals & Planned & research extension\\
Currency conversion and forecasts & Planned & scalability roadmap\\
\bottomrule
\end{tabular}
\end{table}

\chapter{Future work, scalability, and user evaluation}

The next research decisions concern the basket formula, treatment of durable
goods such as cars, missing prices, substitutions, category taxonomy, and the
definition of uncertainty. A versioned category configuration and immutable
monthly input snapshots will support reproducibility. Later stages can add a
database while retaining canonical CSV export, then introduce a web UI and
multi-horizon forecasts with backtesting against seasonal-naive baselines.
Scaling is expected to concern many months and many lines for one user rather
than many users; privacy therefore remains local by default.

The prototype is suitable for the conference's structured interaction study.
Participants can complete a short task sequence: understand the problem,
inspect a receipt row, correct a validation error, run a two-month comparison,
and interpret the difference between price and expenditure indices. Evaluation
can collect usability, UX, functionality, accessibility, innovation,
perceived usefulness, technical robustness, and educational or scientific
impact. Results should be aggregated and reported as user feedback, with
participant privacy preserved and without presenting subjective ratings as
scientific validation of the inflation formula.

\chapter{Conclusion}

PEIT demonstrates a transparent route from manually entered receipts to two
distinct personal economic indicators. Its principal value at this stage is a
reviewable foundation: observations remain traceable, analytical choices are
visible, and the difference between price change and spending behaviour is
preserved. The planned extensions will test whether richer matching, weights,
uncertainty, and forecasts improve usefulness without weakening that evidence
boundary.

\end{document}
